%!TEX root = ../memoria.tex
% =============================================================================
%  Resumen
% =============================================================================

La transformación digital de las últimas décadas ha redefinido profundamente las
formas de comunicación humana, pero también ha dejado al margen a un colectivo
amplio para el que muchas de estas herramientas no fueron concebidas: las
personas mayores. El resultado es una brecha digital que, lejos de reducirse, se
acentúa con la complejidad creciente de las aplicaciones de uso cotidiano y
contribuye al aislamiento de un grupo de población que ya presenta tasas
elevadas de soledad no deseada.

Este Trabajo Fin de Grado presenta el diseño, desarrollo y evaluación de
\emph{ConectaMayor}, una aplicación web accesible cuyo propósito es facilitar
la comunicación entre las personas mayores y sus familiares en un espacio
digital privado, sencillo y seguro. La aplicación se organiza en cinco módulos
—agenda de recordatorios, listado de contactos, mensajería uno a uno, galería
de fotografías compartidas y vinculación familiar mediante código— y contempla
tres roles diferenciados, con permisos adaptados al uso que cada miembro de la
familia realiza del sistema: la persona mayor, el familiar editor y el familiar
de solo lectura.

El sistema se ha implementado con Django, SQLite, Bootstrap 5 y JavaScript
estándar con sondeo AJAX. Esta elección responde a la necesidad de mantener una
base de código manejable en un desarrollo individual y a la adecuación de estas
herramientas a los requisitos del proyecto. El diseño de la interfaz se ha
guiado por las pautas de accesibilidad WCAG 2.1 en su nivel AA y por los
principios del diseño centrado en el usuario, con especial atención a los
aspectos visuales —tipografía, contraste y áreas de interacción— y a la
simplificación del lenguaje empleado.

La validación se ha llevado a cabo mediante una evaluación de usabilidad con
dos personas mayores, siguiendo la metodología propuesta por la Asociación
Interacción Persona-Ordenador (AIPO). Los resultados muestran que la aplicación
cumple los objetivos de sencillez y comprensibilidad fijados, al tiempo que han
permitido identificar puntos de mejora concretos, incorporados al diseño antes
de la entrega final.